
\subsection{The quantum disk}

In this subsection we look at the structure of the quantum disk.
We review the tools and the relevant theorems that will be a~motivation for the subsequent discussion of the quantum
pair of pants.

Consider the closed unit disk ${\mathbb D}=\{\zeta\in{\mathbb C}: |\zeta|\leq 1\}$.
We can represent any holomorphic function inside the disk as a~convergent power series
\begin{gather*}
f(\zeta)=\sum\limits_{n=0}^\infty e_n\zeta^n.
\end{gather*}
The Hardy space on the disk is def\/ined as
\begin{gather*}
H^2({\mathbb D})=\left\{f(\zeta)=\sum\limits_{n=0}^\infty e_n\zeta^n: \sum\limits_{n=0}^\infty|e_n|^2
<\infty\right\}.
\end{gather*}
We def\/ine the multiplication operator by the complex coordinate, $z:H^2({\mathbb D}) \to H^2({\mathbb D})$ by the
formula $f(\zeta)\mapsto \zeta f(\zeta)$.
If $E_n=\zeta^n$ is the orthonormal basis on $H^2({\mathbb D})$, then applying~$z$ to the basis elements produces
$zE_n=E_{n+1}$ for all $n \geq 0$, i.e.,~$z$ is the unilateral shift; moreover, we have the following formula for the adjoint operator to~$z$
\begin{gather*}
z^*E_n=
\begin{cases}
E_{n-1} &\text{for} \quad  n\ge 1,
\\
0  &\text{for} \quad n=0.
\end{cases}
\end{gather*}

Now we consider the $C^*$-algebra generated by~$z$.
This well-known algebra is called the Toeplitz algebra, denoted by $\mathcal{T}$, and has also been termed the quantum
(noncommutative) disk.
This is (partially) based on the following standard results collected here with sketches of proofs which serve as
a~guideline for considerations in the next section.

\begin{Theorem}
The norm of~$z$ is~$1$.
The spectrum of~$z$ is all of ${\mathbb D}$, i.e.,~$\sigma(z)={\mathbb D}$.
\end{Theorem}

\begin{proof}
The norm computation is straightforward.
By the norm calculation it then follows that the spectrum is a~closed subset of the unit disk.
To illustrate that any~$\lambda$ in the interior of ${\mathbb D}$ is an eigenvalue of $z^*$, take $f_\lambda(\zeta) =
\sum\limits_{n=0}^\infty \lambda^n \zeta^n$ and so
\begin{gather*}
z^*f_\lambda(\zeta)=\sum\limits_{n=1}^\infty \lambda^n \zeta^{n-1}=\lambda\sum\limits_{n=1}^\infty
\lambda^{n-1}\zeta^{n-1}=\lambda\sum\limits_{n=0}^\infty \lambda^n\zeta^n=\lambda f_\lambda(\zeta).\tag*{\qed}
\end{gather*}
\renewcommand{\qed}{}
\end{proof}

Let $\mathcal{K}$ be the algebra of compact operators in $H^2({\mathbb D})$.
The next observation tells us how the commutator ideal of $\mathcal{T}$, and $\mathcal{K}$ are related.

\begin{Theorem}
The commutator ideal of $\mathcal{T}$ is the ideal of compact operators.
\end{Theorem}

\begin{proof}
Since $\mathcal{T}$ is generated by~$z$ and $z^*$, the commutator ideal of $\mathcal{T}$ is equal to the ideal generated
by the commutator $[z^*,z]$.
Note that $[z^*,z]=P_{E_0}$, the orthogonal projection onto the span of $E_0$.
Since this one-dimensional projection is a~compact operator, it follows that the commutator ideal of $\mathcal{T}$ is
contained in $\mathcal{K}$.
To prove the opposite inclusion we look at the following rank one operators: $E_{ij}(f)=\langle f, E_i\rangle E_j$.
Notice that $E_{ij}=z^jP_{E_0}(z^*)^i$, hence those operators belong to the commutator ideal of $\mathcal{T}$.
But every compact operator is a~norm limit of f\/inite rank operators, which in turn are f\/inite linear combinations of
$E_{ij}$'s.
This verif\/ies that the commutator ideal of $\mathcal{T}$ contains $\mathcal{K}$.
\end{proof}

In order to state the next result, f\/irst we introduce some more notation.
We identify $H^2({\mathbb D})$, the Hardy space on the unit disk, with the subspace of $L^2(S^1)$ spanned~by
$\{e^{inx}\}_{n\ge0}$.
Also given a~continuous function~$f$ on the unit circle, we denote the multiplication operator by~$f$ as~$M_f$.
Let $P:L^2(S^1)\to H^2({\mathbb D})$ be the orthogonal projection onto $\spn\{e^{inx}\}_{n\ge0}$, then def\/ine
the operator $T_f: H^2({\mathbb D}) \to H^2({\mathbb D})$ by~$T_f=PM_f$.
The operator~$T_f$ is known as a~Toeplitz operator.
Since $\|M_f\|=\|f\|_{\infty}$, $\|T_f\|\le \|f\|_{\infty}$ and hence it is bounded.
We have:

\begin{Theorem}
\label{toeplitz_disk}
The quotient $\mathcal{T}/\mathcal{K}$ is isomorphic to $C(S^1)$, the space of continuous functions on the unit circle.
\end{Theorem}

\begin{proof}
The usual proof constructs an isomorphism between the two algebras.
Notice that for a~continuous function~$f$, we have $T_f\in\mathcal{T}$ and since $T_{e^{ix}}$ is the unilateral shift,
$T_{e^{-ix}}=T_{e^{ix}}^*$.
By the Stone--Weierstrass theorem, every continuous function can be approximated by trigonometric polynomials.
Consequently we can def\/ine a~map $\theta: C(S^1) \to \mathcal{T}/\mathcal{K}$ by $\theta: f \mapsto [T_f]$, the class of
operators $T_f$.

Next we show that $T_f$ is compact if and only if $f\equiv 0$.
Suppose $T_f$ is compact.
Then for a~continuous~$f$ with Fourier series $ \sum\limits_{n=-\infty}^\infty e_n e^{inx}$ we have
\begin{gather*}
T_f\big(e^{ikx}\big)=\sum\limits_{n=0}^\infty e_{n-k}e^{inx}.
\end{gather*}
Thus, the matrix coef\/f\/icients $e_n=(E_{i+n},T_f E_i)$ and since $T_f$ is compact, we must have $(E_{i+n}$,
$T_f E_i)\to0$
as $i\to\infty$ for each f\/ixed~$n$.
Therefore, $e_n=0$ for all~$n$ and hence $f\equiv 0$.
This result means that~$\theta$ is injective.

Next we observe that $T_fT_g-T_{fg}$ is a~compact operator for all continuous $f$, $g$.
If $f$, $g$ are trigonometric polynomials then a~direct calculation shows that $T_fT_g-T_{fg}$ is a~f\/inite rank operator.
The general case then follows by appealing to the Stone--Weierstrass theorem.
As a~consequence, the map~$\theta$ above is a~$C^*$-homomorphism.

The range of~$\theta$ is dense since it contains (the classes of) polynomials in~$z$ and $z^*$.
Then by general $C^*$-algebra theory (see~\cite{Conway} for example)~$\theta$ is an isometry hence the range is closed.
This means that $\Ran(\theta)=\mathcal{T}/\mathcal{K}$ and therefore~$\theta$ is a~$*$-isomorphism.
\end{proof}
